(1) 相邻态差 $\Delta E_n=E_{n+1}-E_n=e^2[n+1/2-q_g/e]/C$。因此加一个电子的代价随栅压线性下降，电容能量的正定性由 $C>0$ 保证；此模型不包含单粒子能级间隔。
(2) 简并点满足 $q_g=(n+1/2)e$，即 $V_g=(n+1/2)e/C_g$。零温基态取使 $|ne-q_g|$ 最小的整数 $n$，故 $\langle n\rangle$ 随 $q_g$ 呈整数量子台阶，周期 $\Delta V_g=e/C_g$。
(3) 两态概率 $P_{n+1}=1/[1+\exp(\Delta E_n/(k_BT))]$，$\langle n\rangle=n+P_{n+1}$。求导 $d\langle n\rangle/dq_g=(e/(Ck_BT))P_{n+1}(1-P_{n+1})$，在简并点最大值为 $e/(4Ck_BT)$，温度升高使台阶变钝。
(4) 两态近似要求 $k_BT$ 远小于相邻充电简并点间的能量尺度 $e^2/C$，且隧穿展宽可忽略。台阶特征电荷宽度约 $Ck_BT/e$；温度不改变相邻简并点间距 $e$，仅改变热平滑宽度。 另外，若存在隧穿寿命展宽，还会出现独立于温度的量子展宽。
